% The "after" of the demo: the same page, after five sentences a reader could say.
% Diff this file against broken.tex to see exactly what each sentence cost.
\documentclass[UTF8,11pt]{ctexart}
\usepackage[a4paper,margin=2.5cm]{geometry}
\usepackage{xcolor}
\usepackage{booktabs}
% xurl 让 URL 可以在任意位置断行，否则一个长链接会把版心顶穿
\usepackage{xurl}
% 页眉：先清空再只放页码，横线随之消失
\usepackage{fancyhdr}
\fancyhf{}
\fancyfoot[C]{\thepage}
\renewcommand{\headrulewidth}{0pt}
\pagestyle{fancy}

\begin{document}

\begin{center}
{\LARGE\bfseries 实验结果}
\end{center}

\vspace{10pt}
下表对比了三种模型的准确率。

\begin{center}
\begin{tabular}{lccc}
\toprule
模型 & 参数量 & 准确率 & 训练时间 \\
\midrule
基线 & 1.2M & 91.3\% & 12 min \\
本文 & 3.4M & 94.8\% & 38 min \\
\bottomrule
\end{tabular}
\end{center}

本节的结果表明，图~\ref{fig:curve} 中的曲线在第一阶段上升最快。实验配置与完整日志见 \url{https://example.com/nndl/experiments/2024/round-07/accuracy-vs-epoch-with-regularization} 这个地址。

\begin{figure}[htbp]
\centering
\fcolorbox{black!35}{black!6}{\parbox[c][3.2cm][c]{0.72\linewidth}{\centering\small 训练曲线占位图}}\caption{训练曲线}\label{fig:curve}
\end{figure}

\end{document}
